% DECISIONES (~1 min cada una). Fuentes: EXP-013, EXP-019, EXP-020.

\begin{frame}{Dropout y L2: no ayudan}
  \begin{columns}[c,onlytextwidth]
    \begin{column}{0.4\textwidth}
      \begin{itemize}
        \item \textbf{Hipótesis}: penalizar la complejidad cierra la brecha
        \item \textbf{Evidencia}: train y validación bajan juntas
        \item \textbf{Resultado}: dropout \res{dropout05-dropout}: \res{loss-dropout05}~pp; las intensidades suaves, sin efecto
      \end{itemize}
      \vspace{0.2cm}
      \decisionblock{Sin dropout ni L2: el límite es el sesgo del modelo, no la varianza.}
    \end{column}
    \begin{column}{0.55\textwidth}
      \begin{figure}
        \centering
        \safegraphics[width=\textwidth,height=0.62\textheight,keepaspectratio]{figures/dropout_curve.png}{Curva de validación del dropout}
        \caption{\footnotesize \figcaptionshort{dropout_curve}}
      \end{figure}
    \end{column}
  \end{columns}
  \speakernote{1}{La curva de validación clásica: estamos en el lado del subajuste. Con L2 pasa lo mismo.}
\end{frame}

\begin{frame}{L2 y combinaciones}
  \begin{columns}[c,onlytextwidth]
    \begin{column}{0.4\textwidth}
      \begin{itemize}
        \item \textbf{Hipótesis}: con augmentation, una penalización suave podría sumar
        \item \textbf{Evidencia}: L2 baja train y validación igual que el dropout
        \item \textbf{Resultado}: sobre BN + rotaciones, ninguna combinación mejora
      \end{itemize}
      \vspace{0.2cm}
      \decisionblock{Modelo: BatchNorm + rotaciones, sin dropout ni L2.}
    \end{column}
    \begin{column}{0.55\textwidth}
      \begin{figure}
        \centering
        \safegraphics[width=\textwidth,height=0.62\textheight,keepaspectratio]{figures/l2_curve.png}{Curva de validación de L2}
        \caption{\footnotesize \figcaptionshort{l2_curve}}
      \end{figure}
    \end{column}
  \end{columns}
  \speakernote{1}{Detalle de las combinaciones en el informe (cuadro de cambios sobre BN + rotaciones).}
\end{frame}

\begin{frame}{BatchNorm}
  \begin{columns}[c,onlytextwidth]
    \begin{column}{0.4\textwidth}
      \begin{itemize}
        \item \textbf{Hipótesis}: estabiliza y permite una lr mayor
        \item \textbf{Evidencia}: con la misma lr, empata
        \item \textbf{Resultado}: antes de la activación y lr \res{phase1-lr}: +\res{gain-bn}~pp
      \end{itemize}
      \vspace{0.2cm}
      \decisionblock{BatchNorm antes de la activación; sin él, con augmentation, se pierden \res{loss-no-bn}~pp.}
    \end{column}
    \begin{column}{0.55\textwidth}
      \begin{figure}
        \centering
        \safegraphics[width=\textwidth,height=0.62\textheight,keepaspectratio]{figures/batchnorm.png}{BatchNorm}
        \caption{\footnotesize \figcaptionshort{batchnorm}}
      \end{figure}
    \end{column}
  \end{columns}
  \speakernote{1}{Por sí solo no reduce la brecha (\res{phase1-gap-acc}~pp), pero es lo que más pesa al combinarlo con la augmentation.}
\end{frame}

\begin{frame}{Data augmentation: rotaciones de 90$^\circ$}
  \begin{columns}[c,onlytextwidth]
    \begin{column}{0.4\textwidth}
      \begin{itemize}
        \item \textbf{Hipótesis}: en imágenes aéreas no hay ``arriba''; rotar conserva la etiqueta
        \item \textbf{Evidencia}: train baja de \res{no-aug-train-acc} a \res{phase2-train-acc}; validación sube
        \item \textbf{Resultado}: +\res{gain-rot90}~pp; brillo y contraste empeoran
      \end{itemize}
      \vspace{0.2cm}
      \decisionblock{Rotaciones de 90$^\circ$ (empatan con volteos y D4).}
    \end{column}
    \begin{column}{0.55\textwidth}
      \begin{figure}
        \centering
        \safegraphics[width=\textwidth,height=0.62\textheight,keepaspectratio]{figures/augmentation.png}{Data augmentation}
        \caption{\footnotesize \figcaptionshort{augmentation}}
      \end{figure}
    \end{column}
  \end{columns}
  \speakernote{1}{Codifica una invarianza que la ffNN no tiene. Sobre este modelo, dropout y L2 siguen sin ayudar.}
\end{frame}

\begin{frame}{Brecha entre train y validación}
  \begin{columns}[c,onlytextwidth]
    \begin{column}{0.4\textwidth}
      \begin{itemize}
        \item Entrega 1: brecha de \res{phase0-gap-acc}~pp
        \item Con rotaciones, a igual número de épocas: \res{phase2-gap-acc}~pp
        \item Entrenando más: \res{final-gap-acc}~pp, con mejor validación
      \end{itemize}
      \vspace{0.2cm}
      \keymessage{La regularización mejora la generalización, pero no elimina la brecha.}
    \end{column}
    \begin{column}{0.55\textwidth}
      \begin{figure}
        \centering
        \safegraphics[width=\textwidth,height=0.62\textheight,keepaspectratio]{figures/overfitting_curves.png}{Curvas de train y validación}
        \caption{\footnotesize \figcaptionshort{overfitting_curves}}
      \end{figure}
    \end{column}
  \end{columns}
  \speakernote{1}{La train del modelo regularizado se mide con augmentation durante la época; con augmentation conviene entrenar más.}
\end{frame}
